\begin{textsample}{NEG Dim 1 – ro – Score 2.00 – ro/ro\_ef\_104.txt}  \label{ex:f1_neg_143}
9.10.2 Competências Específicas da Área de Ensino Religioso COMPETÊNCIAS ESPECÍFICAS DE ENSINO RELIGIOSO Conhecer os aspectos estruturantes das diferentes tradições/movimentos religiosos e filosofias de vida, a partir de pressupostos científicos, filosóficos, estéticos e éticos.
Compreender, valorizar e respeitar as manifestações religiosas e filosofias de vida, suas experiências e saberes, em diferentes tempos, espaços e territórios.
Reconhecer e cuidar de si, do outro, da coletividade e da natureza, enquanto expressão de valor da vida.
Conviver com a diversidade de crenças, pensamentos, convicções, modos de ser e viver.
Analisar as relações entre as tradições religiosas e os campos da cultura, da política, da economia, da saúde, da ciência, da tecnologia e do meio ambiente.
Debater, problematizar e posicionar -se frente aos discursos e práticas de intolerância, discriminação e violência de cunho religioso, de modo a assegurar os direitos humanos no constante exercício da cidadania e da cultura de paz.

% matched lemmas: 
\end{textsample}
